\vspace{-2mm}
\section{Method}
\label{sec:method}

In this section, we first formulate Gaussian-based sensorless 3D ultrasound
reconstruction and define the Gaussian scene representation
(Sec.~\ref{sec:method:problem-formulation-and-scene-representation}). Subsequently,
we introduce an \emph{Acquisition-constrained Frame Geometry Model} (AFGM) to
construct coherent frame geometry without measured probe poses
(Sec.~\ref{sec:method:acquisition-constrained-frame-geometry-model}).
Next, we develop a \emph{Physics-aware Local-Global Renderer} (PLGR) for
physically faithful B-mode synthesis
(Sec.~\ref{sec:method:physics-aware-local-global-renderer}).
Furthermore, a \emph{Contribution-guided Gaussian Adaptation Strategy} (CGAS)
is proposed to regulate the Gaussian population for ultrasound reconstruction
(Sec.~\ref{sec:method:contribution-guided-gaussian-adaptation-strategy}).
The overall framework of ULGSplat is illustrated in
Fig.~\ref{fig:methods-pipeline}.

\vspace{-3mm}
\subsection{Problem Formulation and Scene Representation}
\label{sec:method:problem-formulation-and-scene-representation}

We consider Gaussian-based sensorless 3D ultrasound reconstruction from an
ordered sequence of $N$ B-mode frames,
$\mathcal I=\{\mathbf I_i\}_{i=1}^{N}$. With external tracking, the spatial
configuration of a sweep can be specified by the measured probe poses
$\mathcal T^{\mathrm M}
=
\{\mathbf T_i^{p\rightarrow w}\}_{i=1}^{N}$,
where each rigid transform maps the probe coordinate system of frame $i$ to an
externally referenced world coordinate system $w$. However, in the sensorless
setting, measured poses are inaccessible and reconstruction only requires the
relative geometry among frames within a sweep rather than its absolute
placement in the world coordinate system. In this paradigm, we design AFGM to
construct an acquisition-constrained frame-transform sequence
$\mathcal T^{\mathrm A}
=
\{\mathbf T_{i,\mathrm A}^{p\rightarrow c}\}_{i=1}^{N}$
in a canonical coordinate system $c$, providing the consistent relative frame
geometry required for reconstruction.

Following 3D Gaussian Splatting (3DGS)
\cite{kerbl3DGaussianSplatting2023}, we represent the ultrasound scene in the
canonical coordinate system as a set of $G$ anisotropic Gaussian primitives,
$\mathcal G=\{\mathbf G_g\}_{g=1}^{G}$. Each Gaussian is parameterized by its
center $\boldsymbol\mu_g\in\mathbb R^3$, anisotropic scale
$\mathbf s_g\in\mathbb R_+^3$, unit quaternion $\mathbf q_g$, local echo
amplitude $\alpha_g$, and transmission attribute $\tau_g$. Its covariance is
given by
$\boldsymbol\Sigma_g
=
\mathbf R_g\mathbf S_g\mathbf S_g^\top\mathbf R_g^\top$,
where $\mathbf R_g=\mathbf R(\mathbf q_g)$ and
$\mathbf S_g=\operatorname{diag}(\mathbf s_g)$. In contrast to conventional
3DGS, which models optical appearance using opacity and spherical-harmonic
coefficients, we introduce $\alpha_g$ and $\tau_g$ to characterize local echo
formation and acoustic transmission, respectively, yielding a scene
representation tailored to ultrasound image formation.

Given the ordered B-mode sequence $\mathcal I$, the AFGM-derived frame
transforms $\mathcal T^{\mathrm A}$, and an initialized Gaussian scene
$\mathcal G$, the reconstruction task and loss function are formulated as
\begin{equation}
\begin{aligned}
\mathcal G^{*}
={}&
\arg\min_{\mathcal G}
\sum_{i=1}^{N}
\mathcal L
\left(
\widehat{\mathbf I}_i,
\mathbf I_i
\right),
\quad
\widehat{\mathbf I}_i
=
\mathcal R
\left(
\mathcal G,
\mathbf T_{i,\mathrm A}^{p\rightarrow c}
\right),
\\
\mathcal L
\left(
\widehat{\mathbf I}_i,
\mathbf I_i
\right)
&=
(1-\lambda)
\left\|
\widehat{\mathbf I}_i-\mathbf I_i
\right\|_1
+
\lambda
\left[
1-
\operatorname{SSIM}
\left(
\widehat{\mathbf I}_i,
\mathbf I_i
\right)
\right],
\end{aligned}
\label{eq:overall-reconstruction}
\end{equation}
where $\mathcal R$ denotes PLGR, $\mathcal A$ denotes CGAS, and $\lambda=0.2$
balances pixel fidelity and structural similarity. During optimization, CGAS
periodically updates the Gaussian population via
$\mathcal G\leftarrow\mathcal A(\mathcal G)$, which is omitted from
Eq.~\eqref{eq:overall-reconstruction} for clarity.

\vspace{-3mm}
\subsection{Acquisition-constrained Frame Geometry Model}
\label{sec:method:acquisition-constrained-frame-geometry-model}

Existing sensorless ultrasound reconstruction methods often recover
frame-wise probe motion from adjacent B-mode images using learned motion
estimators, which require additional training and computation and may
accumulate errors when pairwise estimates are chained. Reliable motion
inference from ultrasound appearance is also challenging due to limited
inter-slice overlap, repetitive anatomy, view-dependent speckle, and
substantial intensity variation. In contrast, many ultrasound sweep protocols
exhibit strong acquisition regularity, with approximately parallel imaging
planes and translation-dominant progression rather than large frame-to-frame
rotations. Motivated by this observation, we design the
\emph{Acquisition-constrained Frame Geometry Model} (AFGM) in ULGSplat to
replace frame-wise 6-DoF motion estimation with a compact shared geometry
model that exploits frame order and nominal scan progression to directly
construct the relative frame geometry within a sweep.

Specifically, under the translation-dominant acquisition pattern, we take the
probe coordinate system of a reference frame $i_0$ as the canonical coordinate
system $c$ and assume a shared canonical orientation across the sweep. The
inter-frame progression is then modeled by a shared effective spacing $d$
along a predefined signed unit axis $\mathbf u$, yielding the AFGM-derived
frame transform
\begin{equation}
\mathbf T_{i,\mathrm A}^{p\rightarrow c}
=
\begin{bmatrix}
\mathbf I_3 & (i-i_0)d\mathbf u\\
\mathbf 0^\top & 1
\end{bmatrix},
\label{eq:calibration-geometry-single}
\end{equation}
where $d=\rho d_0$ denotes the shared effective inter-frame spacing,
$d_0>0$ is the nominal acquisition spacing, $\rho>0$ globally adjusts the
scan spacing, and $(i-i_0)$ encodes the relative frame order with respect to
the reference frame. 

This construction yields three geometric properties of the resulting frame
geometry.
First, expressing all frames in the same canonical coordinate system provides
a unified spatial reference, with a homogeneous point
$\bar{\mathbf x}^{p_i}$ in frame $i$ mapped as
$\bar{\mathbf x}^{c}
=
\mathbf T_{i,\mathrm A}^{p\rightarrow c}
\bar{\mathbf x}^{p_i}$,
allowing observations from different B-mode frames to be organized directly
within the same Gaussian scene.
Second, the relative transform between any two
frames $i$ and $j$ is analytically determined by
\begin{equation}
\mathbf T_{j\leftarrow i}^{\mathrm A}
=
\left(
\mathbf T_{j,\mathrm A}^{p\rightarrow c}
\right)^{-1}
\mathbf T_{i,\mathrm A}^{p\rightarrow c}
=
\begin{bmatrix}
\mathbf I_3 & (i-j)d\mathbf u\\
\mathbf 0^\top & 1
\end{bmatrix},
\label{eq:afgm-relative-transform}
\end{equation}
which directly guarantees bidirectional consistency,
$\mathbf T_{i\leftarrow j}^{\mathrm A}
=
(\mathbf T_{j\leftarrow i}^{\mathrm A})^{-1}$.
More importantly, the transforms further satisfy composition consistency,
$\mathbf T_{k\leftarrow j}^{\mathrm A}
\mathbf T_{j\leftarrow i}^{\mathrm A}
=
\mathbf T_{k\leftarrow i}^{\mathrm A}$,
for arbitrary frames $i$, $j$, and $k$, making the resulting relative geometry
independent of the transformation path and avoiding error accumulation from
pairwise motion chaining.

Consequently, AFGM replaces independent frame-wise motion estimation with a
shared acquisition model, eliminating the need for appearance-based inference
and pairwise motion chaining, thereby establishing a consistent spatial
organization of the sweep and a reliable geometric basis for subsequent
Gaussian initialization and rendering.

\vspace{-3mm}
\subsection{Physics-Aware Local-Global Renderer}
\label{sec:method:physics-aware-local-global-renderer}

Conventional 3DGS renders an image by projecting 3D Gaussian primitives onto
the image plane and compositing their opacity-weighted appearance in
front-to-back depth order. Nevertheless, B-mode ultrasound follows acoustic
focusing within a finite elevational slice, so each pixel may aggregate echoes
from multiple scatterers within the local acoustic support; meanwhile, echoes
from deeper structures are modulated by tissues encountered earlier along the
propagation path, making the observed intensity jointly dependent on
finite-slice echo formation and preceding-path transmission.

To address this mismatch, we develop a \emph{Physics-aware Local-Global
Renderer} (PLGR) in ULGSplat, comprising a finite-slice local echo pathway and
a path-dependent global transmission pathway. The local pathway models
finite-slice echo formation by accounting for the elevational interaction of
3D Gaussians with the acoustic slice and their probe-plane geometry, whereas
the global pathway models acoustic transmission by tracing the propagation
path and accumulating the effects of causally preceding structures. The two
responses are then composed for physically faithful B-mode synthesis.

\subsubsection{Finite-Slice Local Echo Formation}

To model local echo formation within the finite acoustic slice, we derive the
Gaussian response through two components: finite-slice elevational
conditioning and probe-plane Gaussian marginalization.

We first establish the probe-space Gaussian representation shared by the local
and global pathways. Given either a measured probe pose
$\mathbf T_i^{p\rightarrow w}$ or an AFGM-derived frame transform
$\mathbf T_{i,\mathrm A}^{p\rightarrow c}$, with either transform written as
$[\mathbf R_i\mid\mathbf t_i]$, Gaussian $g$ is transformed from the
corresponding world or canonical coordinate system into probe space as
\begin{equation}
\boldsymbol\mu^{p}_{ig}
=
\mathbf R_i^\top
\left(
\boldsymbol\mu_g-\mathbf t_i
\right),
\qquad
\boldsymbol\Sigma^{p}_{ig}
=
\mathbf R_i^\top
\boldsymbol\Sigma_g
\mathbf R_i,
\label{eq:probe-space-gaussian}
\end{equation}
yielding the probe-space mean
$\boldsymbol\mu^{p}_{ig}$ and covariance
$\boldsymbol\Sigma^{p}_{ig}$, which are also reused by the global transmission
pathway.

To model finite-slice elevational sensitivity, we condition each probe-space
Gaussian with a zero-centered Gaussian slice profile along the calibrated
elevational direction $\mathbf e_{\mathrm{ele}}$, using a shared effective
slice width $\sigma^{\mathrm{ele}}>0$. With
$\boldsymbol\Lambda^{p}_{ig}
=
(\boldsymbol\Sigma^{p}_{ig})^{-1}$,
the resulting Gaussian product yields the conditioned attributes
\begin{equation}
\begin{aligned}
\boldsymbol\Sigma^{e}_{ig}
&=
\left[
\boldsymbol\Lambda^{p}_{ig}
+
\left(\sigma^{\mathrm{ele}}\right)^{-2}
\mathbf e_{\mathrm{ele}}
\mathbf e_{\mathrm{ele}}^\top
\right]^{-1},
\qquad
\boldsymbol\mu^{e}_{ig}
=
\boldsymbol\Sigma^{e}_{ig}
\boldsymbol\Lambda^{p}_{ig}
\boldsymbol\mu^{p}_{ig},
\\
r^{\mathrm{ele}}_{ig}
&=
\frac{
\left(
\mathbf e_{\mathrm{ele}}^\top
\boldsymbol\mu^{p}_{ig}
\right)^2
}{
\mathbf e_{\mathrm{ele}}^\top
\boldsymbol\Sigma^{p}_{ig}
\mathbf e_{\mathrm{ele}}
+
\left(\sigma^{\mathrm{ele}}\right)^2
+
\varepsilon
},
\qquad
\eta_{ig}
=
\exp
\left(
-\frac{1}{2}r^{\mathrm{ele}}_{ig}
\right),
\end{aligned}
\label{eq:elevational-conditioned-attributes}
\end{equation}
where $\boldsymbol\mu^{e}_{ig}$ and
$\boldsymbol\Sigma^{e}_{ig}$ denote the conditioned mean and covariance,
while $r^{\mathrm{ele}}_{ig}$ measures the normalized elevational separation
and $\eta_{ig}$ gives the corresponding slice-overlap factor.

With the conditioned mean $\boldsymbol\mu^{e}_{ig}$ and covariance
$\boldsymbol\Sigma^{e}_{ig}$, we next apply probe-plane Gaussian
marginalization to derive the image-plane Gaussian geometry for local echo
formation. Given the calibrated lateral and axial directions
$\mathbf e_{\mathrm{lat}}$ and $\mathbf e_{\mathrm{ax}}$, we define the
projection matrix
$\mathbf P_{\mathrm{la}}
=
[\mathbf e_{\mathrm{lat}},
\mathbf e_{\mathrm{ax}}]^\top$
and obtain
\begin{equation}
\begin{aligned}
\boldsymbol\mu^{2D}_{ig}
&=
\mathbf P_{\mathrm{la}}
\boldsymbol\mu^{e}_{ig},
\qquad
\boldsymbol\Sigma^{2D}_{ig}
=
\mathbf P_{\mathrm{la}}
\boldsymbol\Sigma^{e}_{ig}
\mathbf P_{\mathrm{la}}^\top,
\\
\nu_{ig}
&=
\left[
\det
\left(
\boldsymbol\Sigma^{2D}_{ig}
\right)
+
\varepsilon
\right]^{-\gamma},
\end{aligned}
\label{eq:local-image-attributes}
\end{equation}
where $\boldsymbol\mu^{2D}_{ig}$ and
$\boldsymbol\Sigma^{2D}_{ig}$ denote the image-plane mean and covariance,
$\nu_{ig}$ compensates for the Gaussian footprint area, and $\gamma$ controls
the compensation strength.

To avoid exhaustive pixel--Gaussian pairing, we follow the sparse
strategy of 3DGS and construct a local candidate set by jointly considering
the image-space coverage of Gaussian $g$ over pixel $p$ and its local response:
\begin{equation}
\begin{aligned}
\phi_{igp}
&=
\exp
\left[
-\frac{1}{2}
\mathbf d_{igp}^{\top}
\left(
\boldsymbol\Sigma^{2D}_{ig}
\right)^{-1}
\mathbf d_{igp}
\right],
\\
\mathcal N^{\mathrm{loc}}_{ip}
&=
\left\{
g
\;\middle|\;
\boldsymbol\xi^{2D}_{p}
\in
\mathcal B^{2D}_{ig},
\;
\alpha_g\eta_{ig}\nu_{ig}\phi_{igp}
\geq
\tau_{\mathrm{loc}}
\right\}.
\end{aligned}
\label{eq:local-candidate-set}
\end{equation}
Here, $\boldsymbol\xi^{2D}_{p}$ denotes the 2D coordinate of pixel $p$,
$\mathbf d_{igp}
=
\boldsymbol\xi^{2D}_{p}
-
\boldsymbol\mu^{2D}_{ig}$
is its offset from Gaussian $g$, $\phi_{igp}$ denotes the corresponding
image-space response, $\mathcal B^{2D}_{ig}$ is the compact image-space
support, and $\mathcal N^{\mathrm{loc}}_{ip}$ is the resulting candidate set
for local echo accumulation; $\tau_{\mathrm{loc}}$ discards negligible
contributions.

Finally, we jointly combine the echo amplitude $\alpha_g$, slice-overlap factor
$\eta_{ig}$, footprint compensation $\nu_{ig}$, and image-space response
$\phi_{igp}$ to form the full local response in an additive manner, as
ultrasound echo formation differs from the front-to-back alpha compositing
used in optical 3DGS:
\begin{equation}
L_{ip}
=
\sum_{g\in\mathcal N^{\mathrm{loc}}_{ip}}
\alpha_g
\eta_{ig}
\nu_{ig}
\phi_{igp},
\label{eq:local-response}
\end{equation}
where $L_{ip}$ denotes the accumulated local echo response.



\subsubsection{Path-Dependent Global Transmission}

Using the local pathway alone is insufficient to capture ultrasound image
formation, as it ignores preceding structures along the acoustic path and
therefore cannot model path-dependent transmission. To address this limitation,
we introduce path-dependent global transmission through two components:
ray-aware pixel--Gaussian association and causal ray integration.

To identify the Gaussians that can causally affect pixel $p$, we construct a
sparse global candidate set according to the interaction between each
probe-space Gaussian and the corresponding pre-pixel acoustic path. Given the
calibrated pixel coordinate
$\boldsymbol\xi^{2D}_{p}
=
[\xi_{\mathrm{lat},p},\xi_{\mathrm{ax},p}]^\top$,
the acoustic ray in probe space is defined as
\begin{equation}
\mathbf q_{ip}(t)
=
\xi_{\mathrm{lat},p}\mathbf e_{\mathrm{lat}}
+
\left(z_{\mathrm{near}}+t\right)\mathbf e_{\mathrm{ax}},
\qquad
0\leq t\leq t_p,
\label{eq:global-ray}
\end{equation}
where
$t_p=\xi_{\mathrm{ax},p}-z_{\mathrm{near}}$
denotes the propagation distance from the calibrated near plane to the target
pixel. We retain only Gaussians whose compact probe-space support overlaps the
pre-pixel ray corridor and satisfies depth causality:
\begin{equation}
\mathcal N^{\mathrm{glo}}_{ip}
=
\left\{
g
\;\middle|\;
\mathcal B^{\mathrm{glo}}_{ig}
\cap
\mathcal R_{ip}
\neq
\varnothing,\;
\mathbf e_{\mathrm{ax}}^\top
\boldsymbol\mu^{p}_{ig}
\leq
\xi_{\mathrm{ax},p}
\right\},
\label{eq:global-candidate-set}
\end{equation}
where
$\mathcal R_{ip}
=
\{\mathbf q_{ip}(t)\mid0\leq t\leq t_p\}$
denotes the pre-pixel acoustic segment and
$\mathcal B^{\mathrm{glo}}_{ig}$ denotes the compact probe-space support used
for ray association.

For each candidate Gaussian, we quantify its ideal causal overlap with the
pre-pixel acoustic path by first evaluating its pointwise influence along the
ray and then integrating this influence over the complete propagation
interval:
\begin{equation}
\begin{aligned}
\psi_{igp}(t)
&=
\exp
\left[
-\frac{1}{2}
\mathbf d_{igp}(t)^\top
\boldsymbol\Lambda^{p}_{ig}
\mathbf d_{igp}(t)
\right],
\\
\bar{\rho}_{igp}
&=
\int_{0}^{t_p}
\psi_{igp}(t)\,dt.
\end{aligned}
\label{eq:causal-ray-overlap-ideal}
\end{equation}
where
$\mathbf d_{igp}(t)
=
\mathbf q_{ip}(t)-\boldsymbol\mu^{p}_{ig}$
denotes the probe-space offset from Gaussian $g$ to the queried ray point,
$\psi_{igp}(t)$ is its pointwise influence, and $\bar{\rho}_{igp}$ denotes the
ideal ray--Gaussian overlap accumulated over the complete pre-pixel path.

Direct evaluation over the complete pre-pixel interval is computationally
expensive and unnecessary because Gaussian influence rapidly decays away from
its closest interaction with the ray. We therefore approximate the full-path
overlap using a localized integration window, obtained by locating the closest
interaction point $t^{c}_{igp}$ on the infinite acoustic line and clipping its
neighborhood to the causal pre-pixel segment:
\begin{equation}
\begin{aligned}
t_s &= \max(0,t^{c}_{igp}-w),
\qquad
t_e = \min(t_p,t^{c}_{igp}+w),
\\
\rho_{igp}
&=
\int_{t_s}^{t_e}
\psi_{igp}(t)\,dt
\approx
\bar{\rho}_{igp}.
\end{aligned}
\label{eq:localized-ray-overlap}
\end{equation}
where $w$ denotes the integration half-window and $\rho_{igp}$ is the
localized approximation of the ideal ray--Gaussian overlap. The integral is
evaluated using low-order Gauss--Legendre quadrature; if $t_s\geq t_e$, we set
$\rho_{igp}=0$.

Finally, each ray--Gaussian overlap is converted into a transmission factor
through the learnable transmission attribute $\tau_g$, and the preceding
effects are accumulated multiplicatively:
\begin{equation}
\begin{aligned}
T_{igp}
&=
\tau_g
+
\left(
1-\tau_g
\right)
\exp
\left(
-\lambda_{\mathrm{ray}}\rho_{igp}
\right),
\\
T_{ip}
&=
\operatorname{clip}
\left(
\prod_{g\in\mathcal N^{\mathrm{glo}}_{ip}}
T_{igp},
T_{\min},
1
\right),
\end{aligned}
\label{eq:global-transmission}
\end{equation}
where $\lambda_{\mathrm{ray}}$ controls the attenuation strength,
$T_{igp}$ denotes the per-Gaussian transmission factor, and $T_{\min}$ prevents
collapse of the accumulated global transmission $T_{ip}$.

\subsubsection{Local-Global Response Composition}

Finally, we use the global transmission response $T_{ip}$ to modulate the local
echo response $L_{ip}$, with $\beta$ blending $T_{ip}$ with identity
transmission to control the strength of path-dependent modulation:
\begin{equation}
\widehat I_{ip}
=
L_{ip}
\left(
1-\beta+\beta T_{ip}
\right),
\label{eq:local-global-composition}
\end{equation}
where $\beta\in[0,1]$ controls the global transmission strength, with
$\beta=0$ retaining the local response and $\beta=1$ applying the full
transmission response.

By coupling finite-slice local echo formation with path-dependent global
transmission, PLGR preserves local structural detail while incorporating
propagation-dependent context, thereby achieving physically faithful B-mode
synthesis.

\vspace{-4mm}
\subsection{Contribution-guided Gaussian Adaptation Strategy}
\label{sec:method:contribution-guided-gaussian-adaptation-strategy}

Conventional 3DGS density control relies on image-space gradients for
densification and opacity for pruning, which are poorly suited to ultrasound
because speckle can distort gradient cues and B-mode signals are formed by
additive echo accumulation rather than optical alpha compositing. To overcome
these limitations, we design the \emph{Contribution-guided Gaussian Adaptation
Strategy} (CGAS) in ULGSplat, represented by the adaptation operator
$\mathcal A$, with residual-guided densification and path-aware pruning
governed by the candidate sets $\mathcal D$ and $\mathcal P$, respectively.
CGAS allocates additional primitives to under-reconstructed regions and removes
them only when their local and global rendering contributions are jointly weak.

\subsubsection{Residual-Guided Densification}

To identify Gaussians that require additional representation capacity in
under-reconstructed regions, we first define a residual-guided densification
score that measures their contribution to high-error pixels. For an
adaptation-event frame, the pixel residual is
$e_{ip}=|\widehat I_{ip}-I_{ip}|$, while the local contribution of Gaussian
$g$ to pixel $p$ is
$c_{igp}=\alpha_g\eta_{ig}\nu_{ig}\phi_{igp}$. We aggregate these quantities
into the total local contribution
$C_g=\sum_p c_{igp}$ and the residual-weighted contribution
$R_g=\sum_p c_{igp}e_{ip}$, and define the densification score as
\begin{equation}
\omega_g
=
\frac{R_g}{C_g+\varepsilon}
\log(1+C_g),
\label{eq:contribution-aware-densification}
\end{equation}
where a larger $\omega_g$ indicates that Gaussian $g$ contributes more strongly
to under-reconstructed regions and therefore requires additional
representation capacity. We
then construct the densification candidate set by retaining Gaussians with
high positive responsibility scores, subject to a minimum significance
threshold and an approximate per-event budget:
\begin{equation}
\mathcal D
=
\left\{
g
\;\middle|\;
\omega_g
\geq
\max
\left[
Q_{q_{\mathrm d}}
\left(
\{\omega_h\mid\omega_h>0\}
\right),
\tau_{\omega},
\tau_B
\right]
\right\},
\label{eq:densification-selection}
\end{equation}
where $q_{\mathrm d}$ specifies the upper-tail quantile over positive scores,
$\tau_{\omega}$ is the minimum score threshold, and $\tau_B$ is induced by
the per-event budget $B$ to approximately limit the number of selected
Gaussians. Large candidates are refined by retaining the parent and appending
two child Gaussians with reduced raw-scale parameters, whereas smaller
candidates are cloned to increase local representation capacity. The generated
Gaussians inherit the echo and transmission attributes of their parents and
are subsequently refined during optimization.

\subsubsection{Path-Aware Pruning}

To remove Gaussians that are dispensable to both local echo formation and
global transmission, we introduce a global path-contribution measure that
quantifies the accumulated transmission effect of each Gaussian over the
current adaptation interval:
\begin{equation}
A_g
=
\sum_{i\in\mathcal K}
\sum_{p:\,g\in\mathcal N^{\mathrm{glo}}_{ip}}
\left(
1-T_{igp}
\right),
\label{eq:global-path-contribution}
\end{equation}
where $\mathcal K$ denotes the current adaptation interval and $A_g$
accumulates the path-dependent contribution of Gaussian $g$. Together with the
event-frame local contribution $C_g$ and echo amplitude $\alpha_g$, this
measure enables pruning only when a Gaussian is weak in both local and global
rendering:
\begin{equation}
\mathcal P
=
\left\{
g
\;\middle|\;
\alpha_g<\tau_{\alpha},\;
C_g\leq Q_{q_{\mathrm p}}(\{C_g\}),\;
A_g<Q_{q_{\mathrm A}}(\{A_g\})
\right\},
\label{eq:path-aware-pruning}
\end{equation}
where $\tau_{\alpha}$ is the echo-amplitude threshold,
$q_{\mathrm p}$ controls the low local-contribution quantile, and
$q_{\mathrm A}$ controls the path-effect protection level. This criterion
preserves Gaussians with meaningful global path effects even when their local
responses are weak; when global path statistics are unavailable, it reduces
to the corresponding local-aware pruning rule.

Through residual-guided densification and path-aware pruning, CGAS adapts the
Gaussian population to ultrasound-specific rendering demands, expanding
capacity where needed while suppressing redundant primitives.